\begin{center}\textbf{IX. Die Jahre 1827 bis 1843.}\end{center}
\begin{center}\textls[120]{a) \textsc{Abel}s und \textsc{Jacobi}s Entdeckungen. Logarithmisches Potential.}\end{center}
\noindent\hspace*{\parindent}Die Geschichte der Entdeckung der elliptischen Funktionen durch \textsc{Abel}
und \textsc{Jacobi} ist so vielfach und mit Benutzung der sämtlichen zu Gebote
stehenden Quellen so ausführlich dargestellt worden, dass es sich erübrigt,
hier auch nur andeutungsweise auf dieselbe einzugehen. Auch der ausschlaggebende
Einfluss, den die berühmte Stelle im Art. 335 der \textit{Disquisitiones arithmeticae}
sowohl auf \textsc{Abel} als auf \textsc{Jacobi} ausgeübt hat, ist anerkannt; wir können
in bezug auf \textsc{Abel} nur ein gewisses persönliches Moment den bisherigen Darstellungen\footnote{Siehe \textsc{Sylow}, \textit{N. H. Abel, Mémorial etc.} 1902, S. 35, 36.}
hinzufügen, das, wie es scheint, noch nicht bemerkt worden ist.

Es ist nämlich bekannt, dass \textsc{Abel} durch \textsc{Degen} veranlasst worden ist,
sich mit den elliptischen Funktionen zu beschäftigen\footnote{\textit{Mémorial}, Brief von \textsc{Degen} an \textsc{Hansteen}, 21. Mai 1821, \textit{Corresp.} S. 97.}; \textsc{Abel} war Juli—
September 1823 bei \textsc{Degen} in Kopenhagen und beginnt\footnote{Siehe \textsc{Sylow}, a. a. O. S. 27.} 1823 sich mit
Integralrechnung zu beschäftigen. Nun hat \textsc{Schumacher}, wie er in dem oben
erwähnten Briefe an \textsc{Gauss} vom 8. Juni 1816 schreibt\footnote{\textit{Briefwechsel Gauss-Schumacher} I, S. 127.}, \textsc{Degen} mitgeteilt,
\par\clearpage\SourceLeaf{269}\ArithmeticHead{184}{L}{L. SCHLESINGER, ÜBER GAUSS' FUNKTIONENTHEORETISCHE ARBEITEN.}
\noindent dass \textsc{Gauss} sich schon lange mit dem agM. beschäftigt, und es ist sehr wahrscheinlich,
dass \textsc{Degen} von diesem Umstande \textsc{Abel} Mitteilung gemacht hat.
Es könnte also \textsc{Gauss} ausser durch seine Veröffentlichungen, noch durch Vermittlung
von \textsc{Schumacher} und \textsc{Degen} gewissermassen persönlich einen Einfluss
auf \textsc{Abel} ausgeübt haben, wenn dieser Einfluss auch nur darin bestanden hat,
dass er \textsc{Abel} veranlasste, in \textsc{Gauss}' Schriften nach Stellen zu suchen, die sich
auf die Theorie der elliptischen Funktionen beziehen. Dass \textsc{Abel} jene Stelle
in den \textit{Disquisitiones arithm.} gekannt hat, geht z. B. aus dem Briefe an \textsc{Crelle}
vom 4. Dezember 1826\footnote{\textit{Mémorial, Corresp.} S. 52, 53.} hervor; für seine Bekanntschaft mit der \textit{Determinatio
attractionis} haben wir kein direktes Zeugnis. Dagegen hat \textsc{Jacobi}, und zwar
noch ehe er seine eigenen Forschungen auf dem Gebiete der elliptischen
Funktionen begonnen hatte, die \textit{Determinatio attractionis} gründlich studiert,
wie aus den im 2. Bande von \textsc{Crelle}s Journal (1827, S. 227—242) veröffentlichten
Arbeiten: \textit{Über die Hauptachsen der Flächen 2. Ordnung} (\textsc{Jacobi}s Werke
III, S. 45, insbesondere S. 48) und \textit{De singulari quadam duplicis integralis transformatione}
(ebenda S. 55, insbesondere S. 57), hervorgeht.

Von \textsc{Gauss} selbst besitzen wir eine Reihe von brieflichen Äusserungen
über \textsc{Abel} und \textsc{Jacobi}; die eine Briefstelle an \textsc{Bessel} haben wir schon oben
(S. 41) wiedergegeben. Wenn \textsc{Gauss} darin sagt, \textsc{Abel} ȟberhebt mich in
Beziehung auf etwa ein Drittel dieser Sachen der Mühe«, so wird man
dies so aufzufassen haben, dass es sich um den dritten Teil des Materials
handelt, das \textsc{Gauss} für das von ihm geplante grosse Werk über die transzendenten
Funktionen bestimmt hatte (vergl. oben S. 145). Das eine Drittel
wäre die allgemeine Theorie der aus \(F(\alpha,\beta,\gamma,x)\) entspringenden Funktionen,
das zweite die Theorie des agM. und der Modulfunktion, endlich das
dritte, in bezug auf das \textsc{Abel} \textsc{Gauss} zuvorgekommen ist, die Theorie der
elliptischen Funktionen im engeren Sinne. Wenn \textsc{Gauss} noch hinzufügt\footnote{Werke X, 1, S. 248.}:
»Zu meiner Verwunderung erstreckt sich dies [nämlich die Übereinstimmung
seiner Resultate mit denen von \textsc{Abel}] sogar auf die Form und zum Teil auf
die Wahl der Zeichen, so dass manche seiner Formeln wie eine reine Abschrift
der meinigen erscheinen«, so wird man daran zu denken haben, dass \textsc{Abel}
die halben Perioden mit \(\omega,\tilde\omega\) bezeichnet, was in der Tat an \textsc{Gauss}' \(\varpi,\varpi'\) erinnert,
\par\clearpage\SourceLeaf{270}\ArithmeticHead{185}{R}{IXa. ABELS UND JACOBIS ENTDECKUNGEN.}
\noindent und dass infolge dessen die Produktentwicklungen der Zähler und
Nenner bei \textsc{Abel} fast ebenso aussehen wie bei \textsc{Gauss}\footnote{Frappant ist auch z. B. bei \textsc{Abel} (\textsc{Abel}s Oeuvres I, S. 353) die Gleichung
\[
\omega=2\int\frac{dx}{\sqrt{1-x^4}},
\]
die ganz so aussieht wie bei \textsc{Gauss}. Offenbar rührt diese Übereinstimmung daher, dass \textsc{Gauss} sowohl wie
\textsc{Abel} die halbe Periode mit \(\pi\) bezeichnen wollten und dabei die ältere Schreibweise \(\varpi\) dieses Buchstabens
vor Augen hatten; vergl. oben Fussnote \textsuperscript{47)}.}.

Die beiden in den \textit{Astronomischen Nachrichten} von 1827 im Auszug
publizierten Briefe \textsc{Jacobi}s\footnote{\textsc{Jacobi}s Werke I, S. 29 ff.} hat \textsc{Schumacher} an \textsc{Gauss} zur Begutachtung
geschickt. \textsc{Gauss} schickt sie am 4. bezw. 19. August 1827 zurück\footnote{\textit{Briefwechsel Gauss-Schumacher} II, S. 109, 112.} und am
6. August trägt er in das Handbuch 16, Bb (S. 137 ff.) die Werke III, S. 470 ff.
abgedruckte Aufzeichnung ein, die (vergl. Werke III, S. 475, Art. [9]) am
29. August fortgesetzt wird\footnote{Es tritt in dieser Aufzeichnung zu den Charakteristiken \(P,Q,R\) für die drei geraden Thetafunktionen noch die Charakteristik \(S(x,y)\) für die ungerade Thetafunktion hinzu (Werke III, S. 472).}.

In dieser, wie wir wohl sagen können, unter dem unmittelbaren Eindruck
der \textsc{Jacobi}schen Briefe niedergeschriebenen Aufzeichnung finden wir zunächst
Formeln für die Transformation 2. und 3. Ordnung der vier Thetafunktionen.
Wie in der aus demselben Handbuch stammenden Aufzeichnung Werke
X, 1, S. 237 (die auf 1808 zu datieren ist und oben S. 128 bereits besprochen
wurde) wird auch hier (Art. [4], Werke III, S. 472) die allgemeine
Funktion
\[
(\theta,\alpha)=\sum_{k=-\infty}^{+\infty}e^{-\theta(k+\alpha)^2}
=e^{-\theta\alpha^2}\vartheta_{00}\left(\left.\frac{\alpha\theta}{\pi i}\right|e^{-\theta}\right)
\]
\noindent eingeführt, und für sie eine Formel abgeleitet, die die Formeln für die Transformation
zweiter Ordnung der Thetafunktionen zusammenfasst. Setzt man
\(\theta=t\pi\), also \(e^{-\theta}=x,\ 2\theta\alpha=\psi\pi i,\ 2\theta\beta=\chi\pi i\), so schreibt sich diese Formel\begingroup\renewcommand{\thefootnote}{285a)}\footnote{Vergl. \textsc{Jacobi}, 1828, \textsc{Jacobi}s Werke I, S. 257, Gl. (18), (19).}\addtocounter{footnote}{-1}\endgroup
\[
\begin{aligned}
\vartheta_{00}(\psi|x)\vartheta_{00}(\chi|x)
&=\vartheta_{00}(\psi+\chi|x^2)\vartheta_{00}(\psi-\chi|x^2)\\
&\quad+\vartheta_{10}(\psi+\chi|x^2)\vartheta_{10}(\psi-\chi|x^2).
\end{aligned}
\]
\noindent Der folgende Artikel [5] bringt dann einige Rechnungen, die anscheinend darauf
\HistoricalFoot{X 2 Abh. 2.}{24}
\par\clearpage\SourceLeaf{271}\ArithmeticHead{186}{L}{L. SCHLESINGER, ÜBER GAUSS' FUNKTIONENTHEORETISCHE ARBEITEN.}
\noindent hinzielen, von der der Form
\[
\int\frac{du}{\sqrt{1+\mu^2\sin^2u}}
\]
\noindent des elliptischen Integrals entsprechenden Funktion
\[
\operatorname{Sl}\varphi=\frac{q}{r}\frac{\vartheta_{11}(\psi|x)}{\vartheta_{00}(\psi|x)}
\]
\noindent (vergl. oben S. 80) zu einer andern überzugehen, die der in den \textsc{Jacobi}schen
Briefen auftretenden \textsc{Legendre}schen Normalform
\[
\int\frac{dv}{\sqrt{1-k^2\sin^2v}}
\]
\noindent entspricht (vergl. die Fussnote \textsuperscript{115)} oben S. 85).

Es folgt noch (Art. [7]) die Transformation 7. Ordnung und, mit dem
Datum des 29. August versehen, die Transformation 5. Ordnung mit der Modulargleichung.
Bemerkenswert ist die Stelle Art. [10] (Werke III, S. 476),
wo von der allgemeinen Transformation ungerader Ordnung die Rede ist;
\textsc{Gauss} hat also wirklich, wie er an \textsc{Schumacher} schreibt\footnote{\textit{Briefwechsel Gauss-Schumacher} II, S. 112, Brief vom 19. August 1827.}, das in dem zweiten
\textsc{Jacobi}schen Briefe »enthaltene Theorem \mbox{.\,.\,.} ganz leicht aus [seinen] Untersuchungen
über die Transzendenten abgeleitet«. Der auf die Theorie der
Modulfunktion bezügliche Artikel [12] wurde schon oben (Abschn. V c), S. 103,
104) besprochen.

Ein ganz neuer Gedanke tritt hier in überraschender Weise in den Artikeln
[16], [17] auf. Diese beziehen sich auf die \textsc{Laplace}sche partielle Differentialgleichung
\[
\frac{\partial^2 V}{\partial x^2}+\frac{\partial^2 V}{\partial y^2}=0,
\]
\noindent sodass die Vermutung gerechtfertigt erscheint, \textsc{Gauss} habe zu jener Zeit seine
Einsicht in den Zusammenhang zwischen der Theorie der Funktionen komplexer
Variabeln und der des logarithmischen Potentials zu vertiefen gesucht.
Das in dem Artikel [16] ausgesprochene Theorem, dass eine harmonische Funktion,
die auf der Begrenzung eines Gebietes den konstanten Wert \(A\) besitzt,
auch im Innern überall gleich \(A\) ist, bildet die Grundlage für das Prinzip
der analytischen Fortsetzung (vergl. \textsc{Riemann}, Inauguraldissertation, Art. 15,
Werke, 2. Aufl., S. 28); der handschriftliche Nachlass gibt aber keinen Anhaltspunkt
für die Entscheidung der Frage, ob \textsc{Gauss} bis zu diesem Prinzip
\par\clearpage\SourceLeaf{272}\ArithmeticHead{187}{R}{IXb. ALGEBRAISCHE THEOREME.}
\noindent vorgedrungen ist. \textsc{Gauss} beweist dieses Theorem mit Hilfe von Integralsätzen
(ein Beweis, der bekanntlich Bedenken unterliegt) und überträgt es sogleich
auf den Fall des räumlichen Potentials, sodass dieser und der folgende Artikel
zugleich als Vorarbeiten für die 1840 veröffentlichte Abhandlung \textit{Allgemeine
Lehrsätze in Beziehung auf die im verkehrten Verhältnisse des Quadrats
der Entfernung wirkenden Anziehungs- und Abstossungskräfte} (Werke V, S. 195)
gelten können.

